\chapter{Physical Background}
\label{ch:background}

This chapter sets out the analytic theory used throughout the project,
starting from the simplest uniform-field cases and building up to the
dipole field.

\section{Lorentz Force and Basic Assumptions}
\label{sec:lorentz}

For a particle of mass $m$ and charge $q$ with position $\mathbf{r}(t)$
and velocity $\mathbf{v}(t) = \dot{\mathbf{r}}(t)$, the equation of
motion is
\begin{equation}
  m \frac{d\mathbf{v}}{dt} = q\left(\mathbf{E}(\mathbf{r},t) +
    \mathbf{v} \times \mathbf{B}(\mathbf{r},t)\right),
  \label{eq:lorentz}
\end{equation}
where $\mathbf{E}$ and $\mathbf{B}$ are the electric and magnetic fields
\citep{baumjohann1996}. Because the magnetic force is always perpendicular
to $\mathbf{v}$, it does no work on the particle, so in the absence of
an electric field the kinetic energy
$K = \tfrac{1}{2}m|\mathbf{v}|^2$ is conserved. This is used throughout
as an independent check on the numerical integration.

Three simplifying assumptions are made throughout this project.
\begin{enumerate}
  \item \textbf{Non-relativistic motion.} Particle speeds satisfy
    $|\mathbf{v}| \ll c$, so relativistic corrections are negligible.
  \item \textbf{Prescribed fields.} The fields $\mathbf{E}$ and
    $\mathbf{B}$ are given functions of space and time, unaffected by the
    particle's own motion. This holds when particle densities are low
    enough that individual particles do not significantly perturb the
    global field.
  \item \textbf{Single-particle approximation.} Collective effects
    such as collisions, wave--particle interactions and self-consistent
    field generation are neglected. The model describes an isolated
    test particle.
\end{enumerate}

\section{Motion in a Uniform Magnetic Field}
\label{sec:uniform_B}

For a uniform, time-independent field $\mathbf{B} = B_0\hat{\mathbf{z}}$
with $\mathbf{E} = 0$, the velocity can be described by its parallel and
perpendicular components to $\mathbf{B}$:
\begin{equation}
  \mathbf{v} = v_\parallel \hat{\mathbf{z}} + \mathbf{v}_\perp.
  \label{eq:v_decomp}
\end{equation}
The parallel component satisfies $\dot{v}_\parallel = 0$, so motion
along the field is uniform. The perpendicular component rotates at the
gyrofrequency
\begin{equation}
  \Omega = \frac{qB_0}{m},
  \label{eq:gyrofreq}
\end{equation}
tracing a circle of radius
\begin{equation}
  \rho = \frac{v_\perp}{\Omega} = \frac{mv_\perp}{qB_0},
  \label{eq:larmor}
\end{equation}
known as the Larmor radius. Combining uniform parallel motion with
circular gyromotion gives a helical trajectory. The centre of the
gyration, called the guiding centre, stays fixed when the field is
perfectly uniform. A full derivation is given in Appendix~\ref{sec:derivation}.

\section{Motion in Constant Crossed Fields}
\label{sec:exb}

Adding a uniform electric field $\mathbf{E} = E_0\hat{\mathbf{x}}$
perpendicular to $\mathbf{B} = B_0\hat{\mathbf{z}}$ causes the guiding
centre to drift perpendicular to both fields. This is the
$\mathbf{E} \times \mathbf{B}$ drift, with velocity
\begin{equation}
  \mathbf{v}_E = \frac{\mathbf{E} \times \mathbf{B}}{B^2}.
  \label{eq:exb}
\end{equation}
Notably, this drift is independent of particle mass and charge, so it
produces no net current. For the configuration above,
$\mathbf{v}_E = -(E_0/B_0)\hat{\mathbf{y}}$: the guiding centre drifts
in the $-\hat{\mathbf{y}}$ direction while the gyration continues
unchanged.

\section{Gradient-\texorpdfstring{$B$}{B} and Curvature Drifts}
\label{sec:drifts}

Two further drifts arise when the field is spatially non-uniform
\citep{boyd2003}.

\paragraph{Gradient-$B$ drift.}
If $|\mathbf{B}|$ varies across the field, the Larmor radius is larger
where the field is weaker and smaller where it is stronger. The orbit is
no longer a uniform circle and averaging over one gyration leaves a net
drift perpendicular to both $\mathbf{B}$ and $\nabla B$:
\begin{equation}
  \mathbf{v}_{\nabla B} = \frac{m v_\perp^2}{2qB^2}
    \frac{\mathbf{B} \times \nabla B}{B}.
  \label{eq:gradb}
\end{equation}
For the test field
$\mathbf{B} = B_0(1 + \varepsilon x)\hat{\mathbf{z}}$, the gradient
$\nabla B = \varepsilon B_0 \hat{\mathbf{x}}$ gives
$\mathbf{B} \times \nabla B = \varepsilon B_0^2 \hat{\mathbf{y}}$, so
the drift speed reduces to
\begin{equation}
  v_{\nabla B} \approx \frac{m v_\perp^2 \varepsilon}{2qB_0}
  \label{eq:gradb_test}
\end{equation}
in the $\hat{\mathbf{y}}$ direction for positive $q$ and $\varepsilon$.

\paragraph{Curvature drift.}
When field lines are curved, a particle moving along them with parallel
velocity $v_\parallel$ experiences an effective centrifugal force. This
drives a drift perpendicular to both the field and the curvature
direction. For a local radius of curvature $R_c$,
\begin{equation}
  \mathbf{v}_c = \frac{m v_\parallel^2}{qB^2}
    \frac{\mathbf{R}_c \times \mathbf{B}}{R_c^2},
  \label{eq:curv}
\end{equation}
where $\mathbf{R}_c$ points away from the centre of curvature. In a
current-free region ($\nabla \times \mathbf{B} = 0$), the gradient-$B$
and curvature drifts combine into a single expression proportional to
$v_\perp^2 + 2v_\parallel^2$. Both are simultaneously present in the
dipole field studied in Chapter~\ref{ch:results}.

The gradient-$B$ drift can be isolated using a field with straight field
lines (such as $\mathbf{B} = B_0(1+\varepsilon x)\hat{\mathbf{z}}$),
which has a non-zero $\nabla B$ but zero curvature. However, the
curvature drift cannot be tested separately in the same way, because any
physically realisable field with curved field lines necessarily also has
a non-uniform $|\mathbf{B}|$. In the dipole field both drifts act
together and are inseparable.

\section{Magnetic Mirror and Bounce Motion}
\label{sec:mirror}

In a field whose strength increases along the field line, a particle can
be reflected before reaching the strong-field region. The relevant
quantity is the pitch angle $\alpha$, the angle between the velocity
vector and the field direction:
\begin{equation}
  \sin\alpha = \frac{v_\perp}{v},
  \label{eq:pitch}
\end{equation}
so $\alpha = 90^\circ$ means purely perpendicular motion and
$\alpha = 0$ means field-aligned. As the particle moves into stronger
$B$, the magnetic moment
\begin{equation}
  \mu = \frac{m v_\perp^2}{2B},
  \label{eq:mu}
\end{equation}
known as the first adiabatic invariant, is approximately conserved. This
forces $v_\perp^2 \propto B$ to grow, drawing energy from the parallel
motion. The particle reflects --- the mirror point --- when
$v_\parallel = 0$, which occurs when
\begin{equation}
  \sin^2\alpha_\mathrm{eq} = \frac{B_\mathrm{eq}}{B_\mathrm{mirror}},
  \label{eq:mirror}
\end{equation}
where $\alpha_\mathrm{eq}$ is the equatorial pitch angle and
$B_\mathrm{eq}$ the equatorial field strength. Repeated reflection at
both ends produces the oscillatory bounce motion characteristic of
trapped particles.

\section{Dipole Magnetic Field}
\label{sec:dipole}

Earth's geomagnetic field is well approximated near the planet by a
magnetic dipole with moment $M$ aligned along $+\hat{\mathbf{z}}$
\citep{jackson1999}. In Cartesian coordinates,
\begin{equation}
  \mathbf{B}(\mathbf{r}) = \frac{M}{r^5}
    \begin{pmatrix} 3xz \\ 3yz \\ 3z^2 - r^2 \end{pmatrix},
  \label{eq:dipole}
\end{equation}
where $r = |\mathbf{r}|$. This satisfies $\nabla \cdot \mathbf{B} = 0$
and $\nabla \times \mathbf{B} = 0$ everywhere except the origin. On
the dipole axis ($x = y = 0$), the field magnitude reduces to
$|B| = 2M/|z|^3$. The gradient and curvature drifts of
Section~\ref{sec:drifts} are both present throughout.

Field lines are labelled by the $L$-shell parameter, defined so that a
field line crosses the magnetic equatorial plane at distance $LR_E$ from
the centre, where $R_E$ is Earth's radius. The bounce period for a
particle at $L$-shell $L$ with equatorial pitch angle
$\alpha_\mathrm{eq}$ is
\begin{equation}
  \tau_b = \frac{4LR_E}{v} \int_0^{\lambda_m}
    \frac{\cos\lambda\sqrt{1 + 3\sin^2\lambda}}
    {\sqrt{1 - \sin^2\alpha_\mathrm{eq}\,B(\lambda)/B_\mathrm{eq}}}
    \,d\lambda,
  \label{eq:bounce_period}
\end{equation}
where $\lambda_m$ is the mirror latitude from
equation~\eqref{eq:mirror} and
\[
  \frac{B(\lambda)}{B_\mathrm{eq}} =
    \frac{\sqrt{1 + 3\sin^2\lambda}}{\cos^6\lambda}.
\]

\section{The Guiding-Centre Approximation}
\label{sec:gc_theory}

When the gyroradius is small compared to the scale over which the field
varies, it is possible to average over the fast gyromotion and only
track the slow drift of the guiding centre. This is the guiding-centre
(GC) approximation, and it applies when
\begin{equation}
  \frac{\rho}{L_{\nabla B}} \ll 1, \qquad
  \Omega^{-1} \ll T_\mathrm{bounce},
  \label{eq:adiabatic_condition}
\end{equation}
where $L_{\nabla B} = B/|\nabla B|$ is the field gradient scale length.
When these conditions hold, the guiding-centre position
$\mathbf{r}_\mathrm{GC}$ and parallel velocity $v_\parallel$ evolve as
\begin{align}
  \frac{d\mathbf{r}_\mathrm{GC}}{dt} &= v_\parallel \hat{\mathbf{b}}
    + \mathbf{v}_E + \mathbf{v}_{\nabla B} + \mathbf{v}_c,
    \label{eq:gc_r} \\
  \frac{dv_\parallel}{dt} &= -\frac{\mu}{m}\frac{\partial B}{\partial s},
    \label{eq:gc_vpar}
\end{align}
where $\hat{\mathbf{b}} = \mathbf{B}/|\mathbf{B}|$ and $s$ is arc length
along the field line. The second equation is the mirror force: it slows
the particle as it moves into stronger field and reverses it at the
mirror point, directly reproducing the bounce physics of
Section~\ref{sec:mirror}.

Because the system~\eqref{eq:gc_r}--\eqref{eq:gc_vpar} contains no fast
gyromotion, it can be stepped with $\Delta t \gg T_\Omega$, making it far
faster than the full Lorentz solver for the same physical time. A GC
integrator implementing these equations is run alongside the full orbit
solver and compared directly in Section~\ref{sec:gc_vs_full}.
